Fix an integer $k\ge 0$. For a prime $p$, define
\[
F_k(p)\ :=\ \prod_{i=0}^{k}(p^2+i).
\]
Let $P(n)$ denote the largest prime factor of $n$ (with $P(1):=1$).

\noindent\textbf{Question.} Is it true that for every fixed $k$ there exist infinitely many primes $p$ such that
\[
P\bigl(F_k(p)\bigr) = p\ ? 
\]

Equivalently (for $p>k$ so that $p\nmid p^2+i$ for $1\le i\le k$), the condition is:
\[
\max_{1\le i\le k} P(p^2+i)\ \le\ p.
\]
Indeed, the factor $p^2$ ensures that $p$ divides $F_k(p)$, and the requirement that no prime $>p$ divides any of the other factors forces $p$ to be the largest prime divisor.
